\begin{initiativkarte}{Mental Health First Aid (MHFA)}

  % Bild (optional)
  \IfFileExists{bilder/mhfa.jpg}{%
    \begin{center}
      \includegraphics[width=\linewidth,height=5cm,keepaspectratio]{bilder/mhfa.jpg}
    \end{center}
    \vspace{0.4cm}
  }{}

  % Kurzbeschreibung
  \textit{\textcolor{hawgray}{MHFA-Ersthelfende sind geschulte Studierende, die bei psychischen Belastungen unterstützen, zuhören und bei Bedarf an professionelle Hilfe weitervermitteln.}}

  \vspace{0.5cm}
  \hawrule

  % Beschreibung
  \textbf{\textcolor{hawblue}{Was ist MHFA?}}\\[0.3cm]

  Die \textbf{Mental Health First Aid (MHFA)}-Ersthelfenden sind speziell geschulte Studierende der HAW Hamburg.

  Sie bieten erste Unterstützung bei psychischen Belastungen, Krisen und schwierigen Lebenssituationen und helfen dabei, passende Hilfsangebote zu finden.

  \vspace{0.5cm}

  \textbf{Bei welchen Themen kann ich MHFA kontaktieren?}

  Kein Problem ist zu klein. Unterstützung gibt es zum Beispiel bei:

  \begin{itemize}[itemsep=1pt]
    \item Prüfungsangst und Stress
    \item Überforderung im Studium oder Privatleben
    \item Konflikten mit Kommiliton*innen oder Lehrenden
    \item Ängsten und depressiven Symptomen
    \item Privaten Problemen (Familie, Freundschaften, Beziehungen)
    \item Akuten Krisen (z. B. Panikattacken, Suizidgedanken)
  \end{itemize}

  \vspace{0.3cm}

  MHFA kann auch helfen, wenn du professionelle Unterstützung suchst, aber nicht weißt, wo du anfangen sollst.

  \vspace{0.5cm}

  \textbf{Wie läuft der Kontakt ab?}

 Du wählst eine MHFA-Person aus der Liste, schreibst eine E-Mail mit deinem Anliegen (so viel oder wenig wie du möchtest) und vereinbarst anschließend ein vertrauliches Gespräch – online, in der Hochschule oder bei einem Spaziergang.
  
  \vspace{0.5cm}

  % Tags
  \tag{Mental Health}
  \tag{Support}
  \tag{Beratung}
  \tag{Studium}
  \tag{Krise}

  \vspace{0.6cm}
  \hawrule

  % Infozeilen
  \infozeile{\faUsers}{Zielgruppe: Alle Studierenden der HAW Hamburg}
  \infozeile{\faHandsHelping}{Angebot: Erste Hilfe bei psychischen Belastungen}
  \infozeile{\faEnvelope}{Kontakt: über MHFA-Ansprechpersonen (Liste intern)}

\end{initiativkarte}